\ssscp{Receptive to outsiders and outside ideas}{02-03-02-05}
In addition to abundant evidence of Austronesians borrowing
words, metaphors, and even grammatical structures,
but they were also receptive to ``outsiders''
and not ``closed'' in other aspects of life.\footnote{
		The statements here about Austronesian-speaking peoples
		being receptive to outside influence also apply to those
		speaking Papuan languages (e.g. \citealp{Schapper2011-BunakHistory}).}
For example,
\citet{Sahlins1985,Sahlins2008,Fox1995,Grimes2006,Haegerdal2008}
are some of those who have written about the
widespread pattern of accepting ‘outsider kings’
or ‘stranger-kings’ in Austronesian societies.
Given the abundant evidence of people speaking Austronesian languages
borrowing words, metaphors, structures, and being receptive in other ways to outsiders
and their ideas, it is extremely likely that the names for new
flora and fauna encountered as they spread through the target region
were borrowed from other existing terms readily available in
the region, rather than invented from scratch.
The severe under-documentation of languages in the region,
combined with multiple sources of evidence that Papuan languages were once
spoken much more widely in the past (some wiped out by volcanos),
means we may never be able to identify a specific source language.
But that does not make borrowing any less likely.
Arguments in the absence of data do not move us forward.
So we are extrapolating from patterns for which we do have data.
But as shown in \crf{ch:Mar},
there are also Austronesian languages west of the Blust line
which have words of (incrementally) similar form which mean ‘cuscus’,
so a Papuan source is not the only possible explanation.

In response to \citet{Schapper2011}
and \citeauthor{DonohueGrimes2008}' (\citeyear{DonohueGrimes2008}) suggestions that Papuan trade
networks provide a mechanism for the diffusion of words as well
as goods,\footnote{
		Such as those responsible for the human introduction
		of the cuscus from the New Guinea region into western Maluku
		and the Lesser Sundas,
		and bananas, sago, taro and sugarcane as far as western Flores.}
 \citet[274]{Blust2012b} says,
“Since trade is a universal human phenomenon,
it clearly existed in the Papuan world on a local scale before
AN contact, but there is no evidence that it was pan-archipelagic,
or that any Papuan-speaking group had seafaring capabilities
at all comparable to those of AN speakers when they arrived in
the area with the outrigger canoe complex.”

On the “seafaring capabilities”,
we have already shown that there is growing archaeological evidence
that allowed deep water crossings
and open ocean fishing in the region long pre-dating the arrival
of the Austronesian-speaking peoples
\citep{Spriggs2011,OConnor2015,Bellwood2019ed}.
Numerous scholars have noted that the westward spread from the
New Guinea region of bananas,
sugarcane, taro, sago and breadfruit strongly implies a significant
seafaring capability, again predating the Austronesians
\citep{DenhamDonohue2009,DonohueDenham2009,DonohueDenham2010,Schapper2015,
Klamer2019,GrimesCIP}. There is no need to argue whether or not
it was “comparable” to Austronesian capabilities,
since whatever the technology was,
was clearly adequate for the dispersals we are discussing.
Add to that, the \ili{east}
and west prevailing winds shift back
and forth on a yearly seasonal basis,
dominate the weather, agricultural cycles,
and hunting cycles throughout the region,
and are a focus of all boatmen across the islands,
regardless of size or type of boat,
including dugout canoes without outriggers.

But the kind of dispersal of food crops
and the wide distribution of fairly stable terms like {\#}muku
‘banana’ can -- and apparently did -- happen through simple local
networks of contact, without having to appeal to a “pan-archipelagic” trade network,
or to a single Papuan trade language or lingua franca.
Things for which there is,
as yet, little evidence.
(In fact there is evidence of other long-standing trade networks
that contribute to the complexity in the region,
described in \srf{02-02-01}, but not all of them predate the
arrival of Austronesian-speaking peoples.)

